
\chapter{Conclusion}\label{chap:end}

\section{Summary of the Thesis}
This thesis addressed a practical question in situated human--robot interaction: how a Pepper robot can answer questions about its surroundings from perception rather than from generic language-model knowledge. The central design choice was to keep Pepper as the embodied conversational interface while moving expensive perception, memory, and language processing to a local external server. This split made it possible to combine robot-side interaction, multimodal scene analysis, persistent memory, and grounded dialogue in a single working system.

The resulting architecture is intentionally hybrid. QiChat provides deterministic dialogue flow, turn handling, and access to robot actions, while the server maintains the heavier perception and reasoning pipeline: detection, people metadata, scene graph generation, memory association, and language-model-based response generation. The design therefore respects Pepper's limits without reducing the robot to a passive terminal. Pepper remains the social interface; the server supplies the perceptual and semantic machinery needed for scene-aware dialogue.

\section{Main Findings}
The experiments support the central premise of the thesis: perception becomes more useful for dialogue when it is stored as structured memory rather than passed to the language model as an unstructured prompt alone. Scene graphs are a useful intermediate representation because they preserve object identities, attributes, and relations in a form that can be queried and reused across turns. The results also show, however, that such graphs are not produced reliably by prompting alone. Their quality depends on visual grounding, vocabulary design, and the way contextual evidence is presented to the model.

The strongest Set-of-Mark configuration combined numeric object IDs with class labels and bounding boxes. This setting was the most stable across model families, whereas mask overlays were less consistent. Set-of-Mark was also most reliable when paired with a textual scene description. The caption provides scene-level semantics, and the marked image binds those semantics to concrete detector IDs. The visual mark therefore does not replace language context; it makes that context referentially precise.

The vocabulary experiments point in the same direction. The controlled predicate and attribute space is long-tailed, but more labels did not automatically produce better graphs. Compact semantic and frequency-based vocabularies preserved performance surprisingly well, while random reductions degraded it. The important variable was not only vocabulary size, but whether the remaining labels still formed a coherent visual and conversational ontology. Prompt formatting had a smaller but measurable effect, with different model families preferring either a flat list or a structured predicate--attribute separation.

The interaction measurements were more modest, but they clarify where the architecture is practical. Follow-up questions over existing scene memory were answered within a few seconds, which is usable in spoken interaction. Visual description turns were slower because the system must capture an image, transfer it to the server, run perception and generation, and return the response to Pepper. The main bottleneck was not final answer generation, but visual memory construction and the orchestration around it. For deployment, this argues for answering from existing memory whenever possible and moving expensive updates outside the critical response path as the measured latency in the deployment settings was unsatisfactory.

\section{Limitations and Future Work}
The system is complete enough to demonstrate scene-aware dialogue, but it is not an endpoint. Its main limitations are dependence on external compute, sensitivity to VLM prompting, and the remaining gap between strict scene-graph quality and conversational usefulness. The memory layer is also still scene-centric rather than truly long-term: it does not yet maintain a rich model of user history, intention, or shared context across sessions.

Future work should therefore focus on four directions: more robust memory, stronger relation grounding, broader human--robot evaluation, improving latency. The memory layer should handle uncertainty, persistence, and conflicting observations more explicitly. Relation prediction should be improved through better visual VLM prompting, better vocabularies, or more specialized grounding models. The system should be tested in richer user studies, because offline graph metrics and robot-side latency describe only part of interaction quality. Future research should also be directed at measuring user satisfaction with the system which was not done in this thesis. Finally, the latency could be further reduced through additional optimization of model selection, caching, asynchronous processing or other server-side techniques.

The contribution of this thesis is therefore not a final solution to embodied visual dialogue, but a concrete and tested path toward it. It shows that a legacy social robot can be successfully connected to modern perception and language models without abandoning explicit state, controlled grounding, or inspectable memory. Although the implementation is specific to Pepper, the same architecture can be adapted to other robots with visual perception and speech interaction capabilities. We believe that, with even stronger grounding, this direction can support more reliable situated dialogue in human--robot interaction.
